\scanpage{213}
\uline{Groupe fondamental \unsure{motifique|motivique} discret}

Soit $\unsure{\xi|\zeta}\in T(\mathbb Z)$, on pose
\[
\begin{aligned}
\pi_1(X,\xi)&=\mathrm{Aut}(\xi).\\
             &=G_{\xi}(\mathbb Z),
\end{aligned}
\]

\uline{N.B.} En vertu de Borel\inserted{\unsure{-Harish|-Harsh}}, il n'y a qu'un nb fini de $\xi\in T_{\alpha}(\mathbb Z)$ possibles, pour $\alpha$ donné. D'où un nb fini de \uline{groupes fondamentaux discrets}. Sont ils tous associés à des \deleted{\unsure{deux|choses}} places ??

\[
\begin{aligned}
&\xi\in T_{(\alpha)}(k),\ k\text{ un anneau}\\
&M\text{ motif }\in\mathrm{Ob}\,\unsure{\mathscr M|\mathcal M}_{\alpha}\\
&M_{\xi}\text{ fibre en }\xi,\qquad M_{\xi}=\xi(M)\\
&\pi_1(X,\xi)_{\deleted{\unsure{\alpha|\mathrm{an}|0}}}=\mathrm{Aut}(\xi),\quad
G(X,\xi)=\underline{\mathrm{Aut}}(\xi)\\
&\deleted{\unsure{G_{\xi_{\!k}}|G_{\xi}}}\qquad \pi_1(X,\xi)_{\unsure{\alpha|0|\mathrm{an}}}\simeq G(X,\xi)(k)
\end{aligned}
\]

Comparaison avec $\pi_1(X^{\unsure{\mathrm{top}|\mathrm{an}}},\unsure{\xi|a})$ quand $\xi$ associé à une \uline{place} on a ....

\uline{Ex}. $X$ connexe sur $\mathbb C$, $\xi$ défini par un pt de $X$ à valeurs dans $\mathbb C$. Le foncteur $\xi_{\mathbb C}^{\deleted{\unsure{\mathbb C|a}}}$ a une structure \uline{filtrée}. $\xi$ considéré comme foncteur à valeurs dans les modules sous $\pi_1(X,\xi)$ munis d'une \uline{structure de Hodge} est \unsure{pl|pl.|pl\textsuperscript{t}} fidèle ....
% archive sheet number: 212

